% I/O datasheet for the Chipalooza 2026 LVDS transmitter with PLL.
%
% The pin map here is the one in scripts/gen_top.py (PORTS, PIN_USE, PAD_USE),
% which is also what generates the annotation column down the left of the top
% cell. If the two ever disagree, the generator is right and this file is stale.
%
% Build:  latexmk -pdf io_datasheet.tex
\documentclass[a4paper,10pt]{article}

\usepackage[T1]{fontenc}
\usepackage[utf8]{inputenc}
\usepackage[margin=22mm,top=26mm,bottom=24mm]{geometry}
\usepackage{amsmath}
\usepackage{booktabs}
\usepackage{longtable}
\usepackage{array}
\usepackage{xcolor}
\usepackage{fancyhdr}
\usepackage{siunitx}
\usepackage{enumitem}
\usepackage[hidelinks]{hyperref}

\definecolor{accent}{HTML}{2F5D8A}
\definecolor{rule}{HTML}{C7CDD5}
\definecolor{bad}{HTML}{A83F28}

\sisetup{detect-weight, per-mode=symbol, range-units=single}

\newcommand{\net}[1]{\texttt{#1}}
\newcolumntype{L}[1]{>{\raggedright\arraybackslash}p{#1}}

\pagestyle{fancy}
\fancyhf{}
\lhead{\footnotesize\color{accent}\textbf{sg13cmos5l\_chipalooza\_analog\_project}}
\rhead{\footnotesize I/O datasheet}
\cfoot{\footnotesize\thepage}
\renewcommand{\headrulewidth}{0.4pt}

\title{\vspace{-12mm}\textbf{LVDS Transmitter with Integer/Fractional PLL}\\[2mm]
       \large I/O Datasheet}
\author{Chipalooza 2026 \textperiodcentered\ IHP SG13CMOS5L}
\date{Revision of 1 September 2026}

\begin{document}
\maketitle
\thispagestyle{fancy}

\section{Description}

A PRBS-7 pattern generator and an LVDS output driver, clocked either directly
from the reference or from an on-chip PLL whose feedback ratio is programmable
over the shared digital inputs. The port list is
\texttt{verilog/rtl/user\_project\_wrapper\_4a.v} of
\texttt{RTimothyEdwards/sg13cmos5l\_ocd\_chipalooza}, with every bus expanded
into individual pins. Four dedicated analog pads means wrapper \texttt{\_4a},
and therefore slot \texttt{s1} or \texttt{s16} --- per \texttt{config.txt} the
only two with four.

Signal flow: the reference arrives on \net{analog\_pin[0]}; the PLL multiplies
\net{DIV\_RATIO} and halves the result; \net{lvds\_pattern} selects reference or
PLL clock, generates PRBS-7 and hands a differential pair to \net{lvds\_tx},
which drives two of the analog pads.

\section{Pin summary}

All 54 ports of the wrapper. \emph{Not connected} means the pin exists in the
port list but reaches no logic inside this project.

\begin{longtable}{@{}l l L{86mm}@{}}
\toprule
\textbf{Pin} & \textbf{Dir} & \textbf{Function} \\
\midrule
\endfirsthead
\toprule
\textbf{Pin} & \textbf{Dir} & \textbf{Function} \\
\midrule
\endhead
\multicolumn{3}{@{}l}{\textit{Supplies}}\\
\net{vdd\_3v3}      & inout & \SI{3.3}{\volt} --- driver, pre-driver, output stage \\
\net{vss\_3v3}      & inout & \SI{3.3}{\volt} return \\
\net{vdd\_1v2}      & inout & \SI{1.2}{\volt} --- pattern generator, PLL \\
\net{vss\_1v2}      & inout & \SI{1.2}{\volt} return \\
\net{vssio}         & inout & not connected \\
\addlinespace
\multicolumn{3}{@{}l}{\textit{Clock and enable}}\\
\net{clk}           & in    & not connected \\
\net{enable}        & in    & not connected --- see \S\ref{sec:limits} \\
\addlinespace
\multicolumn{3}{@{}l}{\textit{Pattern generator control}}\\
\net{dig\_in[0]}    & in    & \net{clk\_src}: 0 takes the pad, 1 takes the PLL \\
\net{dig\_in[1]}    & in    & \net{en} --- gates the pattern clock \\
\net{dig\_in[2]}    & in    & \net{reset}, active high, seeds the PRBS \\
\addlinespace[2pt]
\multicolumn{3}{@{}p{0.95\linewidth}}{\footnotesize Neither reaches the two
output flops: they run on an ungated clock with \net{RESET\_B} tied high, so
\net{D\_p} and \net{D\_n} are complementary from the first edge after
power-up and the driver settles once rather than at every enable.}\\
\net{dig\_in[3]}    & in    & \net{mode}: 0 clock passthrough, 1 PRBS-7 \\
\net{dig\_in[4]}    & in    & not connected \\
\addlinespace
\multicolumn{3}{@{}l}{\textit{PLL control}}\\
\net{dig\_in[5]}    & in    & PLL \net{ENABLE} \\
\net{dig\_in[6]}    & in    & PLL \net{RESET\_N}, active low \\
\net{dig\_in[16:7]} & in    & \net{DIV\_RATIO[9:0]}, unsigned Q7.3 --- see \S\ref{sec:pll} \\
\net{dig\_in[18:17]}& in    & \net{TEST\_DIV[1:0]} --- selects the \net{TEST\_CLK} divider \\
\net{dig\_in[23:19]}& in    & not connected \\
\addlinespace
\multicolumn{3}{@{}l}{\textit{Digital outputs}}\\
\net{dig\_out[11:0]}& out   & not connected \\
\addlinespace
\multicolumn{3}{@{}l}{\textit{Dedicated analog pads}}\\
\net{analog\_pin[0]}& inout & \net{ref\_clk} --- PLL reference, \SIrange{25}{250}{\mega\hertz} verified \\
\net{analog\_pin[1]}& inout & \net{TEST\_CLK} brought off chip \\
\net{analog\_pin[2]}& inout & LVDS output, non-inverting \\
\net{analog\_pin[3]}& inout & LVDS output, inverting \\
\addlinespace
\multicolumn{3}{@{}l}{\textit{Bias references}}\\
\net{ibias[0]}      & in    & \SI{2}{\micro\ampere} --- pre-driver, mirrored 1:15 internally \\
\net{ibias[1]}      & in    & \SI{2}{\micro\ampere} --- driver, mirrored 1:15 internally \\
\net{vbias}         & in    & not connected \\
\net{analog\_bus[0]}& inout & \SI{2}{\micro\ampere} --- PLL charge pump \\
\net{analog\_bus[1]}& inout & \SI{1.2}{\volt} --- LVDS common-mode reference \\
\net{analog\_bus[3:2]}& inout & not connected \\
\bottomrule
\end{longtable}

\section{Recommended operating conditions}

\begin{center}
\begin{tabular}{@{}l l S S S l@{}}
\toprule
\textbf{Parameter} & \textbf{Pin} & {\textbf{Min}} & {\textbf{Typ}} & {\textbf{Max}} & \textbf{Unit} \\
\midrule
Core supply        & \net{vdd\_1v2}     & 1.08 & 1.20 & 1.32 & V \\
Driver supply      & \net{vdd\_3v3}     & 2.97 & 3.30 & 3.63 & V \\
Reference clock    & \net{analog\_pin[0]} & 25 & 250 & 500 & MHz \\
Bias current       & \net{ibias[1:0]}   & 0.05 & 2    & 10   & \si{\micro\ampere} \\
Charge-pump bias   & \net{analog\_bus[0]} & 0.05 & 2  & 10   & \si{\micro\ampere} \\
Common-mode ref.   & \net{analog\_bus[1]} & {--} & 1.2 & {--} & V \\
\bottomrule
\end{tabular}
\end{center}

\noindent
The bias limits are those of the harness current DACs, whose range is
\SI{50}{\nano\ampere} to \SI{10}{\micro\ampere}. The blocks themselves were characterised at
\SI{30}{\micro\ampere}; the 1:15 pre-mirrors in \net{lvds\_tx} bridge the gap,
so the pins take \SI{2}{\micro\ampere} while the driver and pre-driver still see
the bias they were measured at.

\section{PLL configuration}\label{sec:pll}

The feedback ratio is one unsigned Q7.3 word on \net{dig\_in[16:7]}:
\net{div\_ratio[9:3]} is the integer part and \net{div\_ratio[2:0]} the eighths.
\texttt{fractional\_divider.v} requires the integer part to be at least 4, so
the usable range is \numrange{4.000}{127.875} in steps of \num{0.125}.

\begin{equation*}
  f_{\mathrm{VCO}} = f_{\mathrm{ref}} \times \mathtt{DIV\_RATIO}
  \qquad
  f_{\mathrm{pll\_clk}} = \frac{f_{\mathrm{VCO}}}{2}
  \qquad
  f_{\mathrm{TEST\_CLK}} = \frac{f_{\mathrm{VCO}}}{2^{\,1+\mathtt{TEST\_DIV}}}
\end{equation*}

\noindent
The output divider always halves, so the line rate equals
$f_{\mathrm{ref}} \times \mathtt{DIV\_RATIO} / 2$. The ring oscillator covers
\SIrange{322}{4050}{\mega\hertz} at typical / \SI{1.2}{\volt} /
\SI{27}{\celsius} and does not start below \SI{0.50}{\volt} of control voltage,
so combinations whose VCO frequency falls outside that band will not lock.

\paragraph{Worked example.} \SI{250}{\mega\hertz} on \net{analog\_pin[0]} and a
of \SI{500}{\mega\bit\per\second} wants $\mathtt{DIV\_RATIO} = 4.0$, code
$4.0 \times 8 = 32 = \mathtt{0000100000}_2$. Only bit 5 is set, which is
\net{dig\_in[12]}; every other ratio bit is driven low. The VCO then runs at
\SI{1}{\giga\hertz}.

\paragraph{Fractional operation.} The divider is a plain $N/N{+}1$ accumulator
with no delta-sigma shaping. At a fractional ratio the instantaneous period
alternates and only the average over a full accumulator cycle is the intended
frequency; expect corresponding spurs on the control voltage.

\section{Electrical characteristics}

Simulated at the top level, typical corner, \SI{27}{\celsius}, over
\SI{6}{\micro\second}, with the reference at \SI{250}{\mega\hertz} and
$\mathtt{DIV\_RATIO} = 4.0$. Load is \SI{49.9}{\ohm} + \SI{49.9}{\ohm} across
the two output pads with the centre tap giving $V_{OS}$. \textbf{No pad or ESD
model is included}, so $V_{OS}$ here is not the compliance figure.

\begin{center}
\begin{tabular}{@{}l l S S l@{}}
\toprule
\textbf{Parameter} & \textbf{Symbol} & {\textbf{Measured}} & {\textbf{Limit}} & \textbf{Reference} \\
\midrule
Differential output   & $|V_{OD}|$      & 380.3 & {247\,\ldots\,454} & TIA/EIA-644-A 4.1.1 \\
Offset voltage        & $V_{OS}$        & 1193.9 & {1125\,\ldots\,1375} & 4.1.2 \\
Offset, peak-to-peak  & $V_{OS,pp}$     & 118.1 & {< 150} & 4.1.5 \\
Offset p-p, first \SI{100}{\nano\second} & & 159.4 & {} & cold start, see \S\ref{sec:limits} \\
\addlinespace
PLL output frequency  & $f_{pll\_clk}$  & 500.0117 & {500} & target \\
Frequency error       &                 & 23.3 & {} & ppm \\
Control voltage       & $V_{CTRL}$      & 581.4 & {} & mV, settled \\
\addlinespace
Driver supply current & $I_{3v3}$       & 5.74 & {} & mA \\
Core supply current   & $I_{1v2}$       & 0.341 & {} & mA \\
\bottomrule
\end{tabular}
\end{center}

\noindent
Voltages in \si{\milli\volt} unless noted. The frequency and control-voltage
figures come from the co-simulated PLL, in which the phase detector and both
dividers are the RTL of \texttt{macros/pll\_digital} executed through ngspice's
\texttt{d\_cosim}. Supply currents are from the LVDS bench, which clocks the
pattern generator from the reference directly.

\section{Known limitations}\label{sec:limits}

\begin{itemize}[leftmargin=*]
  \item \textbf{No ESD structures or pad models.} Clause 4.1.4 and 4.1.5 figures
        depend on pad capacitance, which is not represented here.
  \item \textbf{\net{enable} has no load.} The pattern clock is gated by
        \net{dig\_in[1]} alone. The harness masks \net{dig\_in} to zero for an
        unselected project, so an unselected project is still idle --- but
        \net{select} and \net{dig\_ena} high with \net{enable} low leaves the
        block running with the project enable deasserted.
  \item \textbf{$\Delta V_{OS}$ is not measured at the top level.} The
        \SI{50}{\milli\volt} limit of clause 4.1.2 is checked only in the macro's
        own DC bench.
  \item \textbf{Offset margin is asymmetric, and the reference is unsettled.}
        The benches drive \net{analog\_bus[1]} at \SI{1.2}{\volt}, where
        $V_{OS}$ sits \SI{69}{\milli\volt} above the lower limit of clause
        4.1.2 and \SI{181}{\milli\volt} below the upper. The rest of the
        documentation calls this reference \SI{1.25}{\volt}, which is where the
        clause centres and would split the margin evenly. Which of the two the
        harness bias generator can actually deliver is unverified, so the
        figures here are the \SI{1.2}{\volt} ones and the discrepancy is open.
  \item \textbf{$V_{OS,pp}$ is switching feedthrough, not settling.} It is a
        \SI{93}{\milli\volt} dip at every single data edge, and it starts in
        the pre-driver rather than the driver. The falling gate edge is
        \SI{103.5}{\pico\second} against \SI{112.0}{\pico\second} rising, so at
        each crossing the gate common mode $(V_{In_p}+V_{In_n})/2$ dips
        \SI{78}{\milli\volt}; that lands on the output common mode at a ratio
        of 1.11, identically in both data directions. The cross-coupled
        neutralisation caps cannot cancel it --- they remove differential
        feedthrough and this term is common mode --- and the CMFB loop sees
        only \SI{36.8}{\milli\volt} of it through its sense filter, far too
        slow for a \SI{165}{\pico\second} event in any case. The tail current
        imbalance was measured and ruled out: \SI{3.3}{\femto\coulomb} where
        \SI{34.6}{\femto\coulomb} would be needed. Two levers are characterised
        but not adopted --- the pre-driver's last PMOS has a narrow optimum at
        \SIrange{36}{38}{\micro\metre} giving \SI{102}{\milli\volt}, and a
        common-mode RC damper across the outputs gives \SI{97}{\milli\volt} for
        \SI{2.6}{\percent} of $|V_{OD}|$.
  \item \textbf{The first \SI{100}{\nano\second} are outside the compliance
        window.} $V_{OS,pp}$ reaches \SI{159.4}{\milli\volt} there while the
        driver's common-mode loop settles from the DC operating point. The
        benches measure the compliance figures from \SI{150}{\nano\second}
        onwards. This is what the ungated, reset-free output flops improved:
        the same window measured \SI{288.3}{\milli\volt} when the pair was not
        complementary until the pattern was enabled.
  \item \textbf{\net{analog\_pin[1]}} is driven by \net{TEST\_CLK} with no pad
        driver in front of it.
  \item \textbf{\net{core\_p} / \net{core\_n}} cross from \net{vss\_1v2} into the
        pre-driver on \net{vss\_3v3}; the two returns have to be tied in the
        frame.
  \item \textbf{Simulation needs a build step.} The PLL's digital half runs as a
        Verilator-compiled shared object under \texttt{d\_cosim}. Run
        \texttt{make pll-cosim-so} before any top-level simulation; the Makefile
        and the Simulate arrow in xschem both do this automatically.
\end{itemize}

\end{document}
